\documentclass[12pt]{article}
\usepackage{amsmath, amssymb, amsthm}
\usepackage[margin=1in]{geometry}
\usepackage{hyperref}

\title{\vspace{-2em}Proposition 5: Dirac Spectrum Rigidity\\
\large From Paper XXII: The Almost-Quaternionic Mock Modular Correspondence}
\author{}
\date{}

\newtheorem{prop}{Proposition}
\setcounter{prop}{4} % Start at Proposition 5

\begin{document}
\maketitle

\begin{prop}[Second Rigidity Witness: Strict Minimum Inequality]
The squared eigenvalue of the Spin Dirac operator on $(T^2, \tau = i\tau_2)$ with determinant line bundle $L = \mathcal{O}(3)$ attains a minimum of $\pi^2$ at $m=0, n=1$, and is strictly greater than $\pi^2$ for all non-zero winding numbers $m$.
\end{prop}

\begin{proof}
To establish that the squared eigenvalue $\lambda^2(m, n, \tau_2)$ strictly exceeds the absolute minimum $\pi^2$ for any non-zero winding number, let $m, n \in \mathbb{Z}$ and $\tau_2 \in \mathbb{R}$, subject to the physical constraints $m \neq 0$ and $\tau_2 > 0$.

The exact squared eigenvalue is defined as:
\begin{equation}
    \lambda^2(m, n, \tau_2) = \frac{4\pi^2}{\tau_2^2} m^2 + 4\pi^2 \left(n - \frac{3}{2}\right)^2
\end{equation}

Because $n$ is strictly confined to the integers, the continuous distance to the vortex center $\left(n - 3/2\right)^2$ must respect the discrete boundary conditions. Given that $n \le 1$ or $n \ge 2$, we obtain the sharp lower bound:
\begin{equation}
    \left(n - \frac{3}{2}\right)^2 \ge \frac{1}{4}
\end{equation}

Multiplying this by the global scale factor $4\pi^2$ establishes a rigid lower bound for the discrete mode contribution:
\begin{equation}
    4\pi^2 \left(n - \frac{3}{2}\right)^2 \ge \pi^2
\end{equation}

Concurrently, since the winding number is non-vanishing ($m \neq 0$) and the modular parameter volume is strictly positive ($\tau_2 > 0$), the kinetic winding contribution must be strictly positive:
\begin{equation}
    \frac{4\pi^2}{\tau_2^2} m^2 > 0
\end{equation}

Combining these independent inequalities by direct linear addition yields the strict inequality for the total spectrum:
\begin{equation}
    \frac{4\pi^2}{\tau_2^2} m^2 + 4\pi^2 \left(n - \frac{3}{2}\right)^2 > \pi^2
\end{equation}

This completes the verification that any state with a non-zero winding number occupies a strictly higher energy level than the $m=0, n=1$ ground state.
\end{proof}

\vspace{2em}
\noindent\textbf{Code Availability \& Mechanical Verification}\\
\rule{\textwidth}{0.4pt}
The central algebraic equivalences and inequalities of this proposition have been formally verified using the \textbf{Lean 4} interactive theorem prover. The complete source code, geometric definitions, and compilation instructions demonstrating the absolute absence of logical gaps in the mode scaling are permanently available in the corresponding GitHub repository under \texttt{Paper\_XXII\_Letter/Prop5\_Dirac\_Spectrum.lean}.

\end{document}
